\documentclass[
  spanish,
  letterpaper, oneside
]{article}

\def\documenttitle {Cuestionario diagnóstico sobre interaprendizaje y colaboración digital}
\def\documentsubtitle {Diagnóstico inicial}
\def\documentsubject {Actividad 1 - Interaprendizaje en ambientes virtuales}

\def\documentauthor {Martin Jonathan de la Cruz}
\def\coursename {Interaprendizaje en ambientes virtuales}
\def\coursecode {004}

\def\universityname {Universidad Abierta y a Distancia de México}
\def\universityfaculty {Licenciatura en Derecho}
\def\universitydepartment {Interaprendizaje en ambientes virtuales}
\def\universitydepartmentimage {departamentos/UnADM}
\def\universitydepartmentimagecfg {height=1.57cm}
\def\universitylocation {Roma Norte, Ciudad de México}

\def\authortable {
  \begin{tabular}{ll}
    \textbf{Alumno:}
    & \begin{tabular}[t]{l}
      Martin Jonathan de la Cruz \\
    \end{tabular} \\
    Matrícula: & ES2611202040 \\
    Figura docente: & Janeth Santana Ramírez \\
    Semestre/Bloque: & 2 / 1 \\
    Tipo/Créditos: & Optativa / 6 \\
    & \\
    \multicolumn{2}{l}{Fecha de realización: \today} \\
    \multicolumn{2}{l}{\universitylocation}
  \end{tabular}
}

\input{template}
\usepackage{helvet}
\renewcommand{\familydefault}{\sfdefault}
\usepackage{array}
\usepackage{longtable}
\usepackage{pdflscape}
\usepackage{tikz}
\usetikzlibrary{arrows.meta,positioning,calc,matrix,fit,shapes.geometric,shadows.blur}
\setcitestyle{authoryear,open={(},close={)}}

\def\coverwatermarkenabled {true}
\def\coverwatermarkimage {img/departamentos/UnADM.pdf}
\def\coverwatermarkopacity {0.16}
\def\coverwatermarkwidth {0.72\paperwidth}
\def\coverwatermarkxshift {0cm}
\def\coverwatermarkyshift {-0.35cm}

\newcommand{\insertcoverwatermark}{%
  \ifthenelse{\equal{\coverwatermarkenabled}{true}}{%
    \AddToShipoutPictureBG*{%
      \begin{tikzpicture}[remember picture,overlay]
        \node[
          opacity=\coverwatermarkopacity,
          inner sep=0pt,
          xshift=\coverwatermarkxshift,
          yshift=\coverwatermarkyshift
        ] at (current page.center) {%
          \includegraphics[width=\coverwatermarkwidth]{\coverwatermarkimage}%
        };
      \end{tikzpicture}%
    }%
  }{}%
}

\begin{document}

\insertcoverwatermark
\templatePortrait
\templatePagecfg
\onehalfspacing

\begin{abstractd}
  El documento responde un cuestionario diagnóstico sobre interaprendizaje,
  colaboración digital y diversidad cultural y lingüística. Las notas de clase se
  emplean como pivote para delimitar preguntas y enfoque; las referencias
  académicas y normativas se sostienen en fuentes sólidas sobre aprendizaje
  colaborativo, competencias digitales, interculturalidad, derechos lingüísticos,
  acceso a la justicia y uso responsable de inteligencia artificial en educación.
\end{abstractd}

\templateIndex
\templateFinalcfg

\section{Introducción}

El cuestionario diagnóstico permite fijar una línea base: qué se entiende por
interaprendizaje, cómo se percibe la capacidad para colaborar a distancia, qué
herramientas digitales se usan y de qué modo la diversidad cultural y lingüística
se conecta con la formación jurídica. En lugar de citar las notas internas como
fuente final, se usan como guía de organización, mientras las ideas se contrastan
con literatura académica, fuentes institucionales y marco jurídico vigente.

El interaprendizaje se relaciona con una tradición pedagógica que entiende el
aprendizaje como proceso social: las personas construyen conocimiento al dialogar,
contrastar perspectivas y resolver tareas con otras personas. Esta idea es
compatible con el enfoque sociocultural de \citet{vygotsky1978mind} y con la
literatura de aprendizaje colaborativo que subraya interdependencia positiva,
responsabilidad individual, interacción promotora y evaluación del proceso grupal
\citep{johnsonJohnson1999learning}. En la formación jurídica, esto importa porque
el Derecho se aprende y se ejerce mediante argumentación, escucha, revisión de
supuestos y deliberación razonada.

\section{Cuestionario sobre interaprendizaje y colaboración digital}

El cuestionario aborda el punto de partida para aprender con otras personas en
entornos virtuales. Sus reactivos exploran la comprensión inicial del
interaprendizaje, las condiciones del trabajo a distancia, la seguridad en el uso
de herramientas digitales, la experiencia en autoevaluación y coevaluación, y la
importancia de la diversidad cultural y lingüística para una formación jurídica
orientada al acceso a la justicia. En conjunto, las respuestas muestran una base
favorable para colaborar, pero también la necesidad de convertir la disposición
inicial en acuerdos, evidencias y criterios verificables.

\begingroup
\scriptsize
\setlength{\tabcolsep}{3.5pt}
\renewcommand{\arraystretch}{1.25}
\begin{longtable}{@{}>{\raggedright\arraybackslash}p{0.06\textwidth}>{\raggedright\arraybackslash}p{0.30\textwidth}>{\raggedright\arraybackslash}p{0.27\textwidth}>{\raggedright\arraybackslash}p{0.29\textwidth}@{}}
\caption{Cuestionario diagnóstico sobre interaprendizaje y colaboración digital}\label{tab:cuestionario-interaprendizaje}\\
\toprule
\textbf{No.} & \textbf{Pregunta} & \textbf{Respuesta} & \textbf{Justificación} \\
\midrule
\endfirsthead
\toprule
\textbf{No.} & \textbf{Pregunta} & \textbf{Respuesta} & \textbf{Justificación} \\
\midrule
\endhead
\bottomrule
\endfoot
\bottomrule
\endlastfoot
\midrule
1 & Cuando escuchas la palabra ``interaprendizaje'', ¿con qué idea la relacionas más en este momento? & Con aprender a partir del intercambio de ideas, experiencias y perspectivas con otras personas. & El interaprendizaje se centra en aprender mediante intercambio con otras personas; esta noción coincide con enfoques socioculturales y colaborativos del aprendizaje \citep{vygotsky1978mind,johnsonJohnson1999learning}. \\
\midrule
2 & En tu experiencia u opinión, ¿qué condiciones ayudan a que un trabajo en equipo a distancia funcione adecuadamente? & Comunicación clara y constante, roles definidos, acuerdos y plazos claros, herramientas digitales adecuadas, responsabilidad, respeto y seguimiento continuo de tareas. & La colaboración remota requiere reglas explícitas para evitar ambigüedad, pérdida de versiones y falta de trazabilidad. \\
\midrule
3 & ¿Cómo valorarías tu habilidad actual para trabajar en equipo a distancia? & Media: se puede colaborar, aunque conviene mejorar organización, comunicación o seguimiento. & Es una respuesta de autopercepción; por eso debe verificarse después con evidencias como cumplimiento de plazos, calidad de aportes y participación en revisiones. \\
\midrule
4 & ¿Qué entiendes por diversidad cultural y lingüística? & Coexistencia de distintas culturas, lenguas, costumbres y formas de ver el mundo dentro de una sociedad. & La diversidad no es un dato decorativo: afecta comprensión, participación, comunicación y acceso efectivo a derechos. \\
\midrule
5 & En la formación jurídica, ¿para qué puede servir reconocer la diversidad cultural y lingüística? & Para garantizar acceso a la justicia, respetar derechos de todos los grupos, aplicar el Derecho con enfoque intercultural y evitar discriminación. & En México, la Constitución reconoce la composición pluricultural y prohíbe la discriminación; además, la legislación protege derechos lingüísticos de los pueblos indígenas \citep{cpeum2026,leyDerechosLinguisticos2024}. \\
\midrule
6 & ¿Qué herramientas digitales has utilizado para estudiar, comunicarte o trabajar en equipo? & Chats o grupos de mensajería; videollamadas; documentos compartidos; plataformas educativas; herramientas de inteligencia artificial; bibliotecas digitales, bases de datos o repositorios académicos. & Cada herramienta requiere una función clara: coordinar, deliberar, redactar, entregar, apoyar ideas o verificar fuentes. Las competencias digitales implican uso crítico, seguro y responsable \citep{unescoDigitalLiteracy2018}. \\
\midrule
7 & ¿Qué tan seguro/a te sientes usando herramientas digitales para aprender y colaborar? & Seguro/a: se utilizan con frecuencia, aunque a veces se necesita apoyo. & La seguridad digital inicial es útil, pero no equivale a competencia avanzada; debe complementarse con trazabilidad, protección de información y revisión crítica. \\
\midrule
8 & ¿Has participado antes en actividades de autoevaluación? & Sí, varias veces. & La autoevaluación favorece reconocimiento del propio desempeño, fortalezas y áreas de mejora. \\
\midrule
9 & ¿Has participado antes en actividades de coevaluación, es decir, valoración entre compañeras y compañeros? & Sí, pocas veces. & La coevaluación requiere criterios compartidos y evidencia observable para evitar juicios arbitrarios o sesgos de afinidad. \\
\midrule
10 & ¿Qué esperas aprender, practicar o fortalecer en esta asignatura? & Fortalecer trabajo colaborativo a distancia, uso de herramientas digitales, enfoque intercultural y habilidades de comunicación aplicadas al ámbito jurídico. & La expectativa integra los ejes de la asignatura: colaboración, competencia digital, interculturalidad y comunicación jurídica. \\
\end{longtable}
\endgroup

La lectura conjunta del cuestionario muestra una línea base favorable, pero no
conclusiva. Hay comprensión inicial del interaprendizaje, uso funcional de
herramientas digitales y reconocimiento de la diversidad cultural y lingüística;
sin embargo, varias respuestas son autoperceptivas. Por ello, las siguientes
actividades deben convertir esas respuestas en evidencias observables: acuerdos
de equipo, productos colaborativos, fuentes verificables, revisión entre pares y
autoevaluación con criterios explícitos. El uso de inteligencia artificial, cuando
aparezca, debe mantenerse bajo supervisión humana, protección de datos y revisión
crítica de fuentes \citep{unescoGenerativeAI2023}.

% \section{Matriz diagnóstica de partida}
%
% \begin{longtable}{>{\bfseries}p{0.25\textwidth}p{0.29\textwidth}p{0.38\textwidth}}
% \toprule
% Aspecto & Estado inicial & Criterio para siguientes actividades \\
% \midrule
% Interaprendizaje & Se comprende como intercambio de ideas y experiencias. & Exigir contraste de perspectivas antes de cerrar conclusiones. \\
% Trabajo remoto & Requiere comunicación, roles, acuerdos y seguimiento. & Formalizar protocolo de equipo y evidencia de acuerdos. \\
% Competencia digital & Hay uso funcional de varias herramientas. & Asignar una función clara a cada recurso y evitar duplicidades. \\
% Diversidad cultural y lingüística & Se vincula con convivencia, derechos y no discriminación. & Incorporar preguntas de enfoque intercultural en el análisis jurídico. \\
% Autoevaluación y coevaluación & Hay más experiencia individual que entre pares. & Usar rúbricas breves con evidencia observable. \\
% \bottomrule
% \end{longtable}

% ============================================================================
% TRAZABILIDAD EDITORIAL PARA MOTOR INTELIGENTE
% Las notas internas funcionan como pivote de consigna y memoria de aula, no como
% referencias bibliográficas finales. Si el insumo es cuestionario, el desarrollo
% visible debe preservar reactivos, respuestas y justificaciones, preferentemente
% en tabla compacta o lista estructurada. Las citas visibles deben sostenerse en
% fuentes académicas, normativas, institucionales o en línea verificables.
% ============================================================================

\section{Conclusión}

El cuestionario confirma que el interaprendizaje en ambientes virtuales debe
entenderse como una práctica organizada, no como simple participación en línea.
El curso puede lograr su propósito si convierte la disposición a colaborar en
reglas verificables: roles, acuerdos, herramientas con función, respeto a la
diversidad y evaluación basada en evidencia.

En la formación jurídica, esta línea base es especialmente importante. El trabajo
profesional exige argumentar con otras personas, reconocer contextos culturales y
lingüísticos, usar tecnologías con responsabilidad y mantener rigor individual.
Por eso, el cuestionario inicial no es un trámite: es el punto de partida para
pasar de la autopercepción a la evidencia y de la colaboración declarada a un
interaprendizaje jurídicamente útil.

\clearpage
\nocite{unadmSitioWeb,unadmMallaDerecho2024,vygotsky1978mind,johnsonJohnson1999learning,unescoDigitalLiteracy2018,unescoInterculturalCompetences2013,unescoGenerativeAI2023,cpeum2026,leyDerechosLinguisticos2024}
\bibliography{interaprendizaje-en-ambientes-virtuales}

\end{document}